\chapter*{Eidesstattliche Erklärung}
\addcontentsline{toc}{chapter}{Eidesstattliche Erklärung}
% Die aktuelle Version der Eidesstättlichen Erklärung kann beim zuständigen Prüfungsamts gefunden werden.
% Nachfolgend ist die Version für Bachelor und Master des Prüfungsamts Mathe/Informatik verglichen am 22.04.2025.

Hiermit versichere ich, dass die vorliegende Arbeit über \textit{Reengineering verteilter Legacy-Systeme unter Einhaltung von Delivery-Garantien: Entwurf und Evaluation eines asynchronen Proxy-Systems} selbstständig von mir und ohne fremde Hilfe verfasst worden ist, dass keine anderen Quellen und Hilfsmittel als die angegebenen benutzt worden sind und dass die Stellen der Arbeit, die anderen Werken -- auch elektronischen Medien und KI Tools -- dem Wortlaut oder Sinn nach entnommen wurden, auf jeden Fall unter Angabe der Quelle als Entlehnung kenntlich gemacht worden sind. Mir ist bekannt, dass es sich bei einem Plagiat um eine Täuschung handelt, die gemäß der Prüfungsordnung sanktioniert werden kann.

\vspace{1cm}

\noindent \parbox{20em}{\hrulefill}

\noindent Jad Haddad, Münster, \today

\vspace{1cm}

\noindent Ich erkläre hiermit, dass ich Kenntnis von einer zum Zweck der Plagiatskontrolle vorzunehmenden Speicherung der Arbeit in einer Datenbank sowie von ihrem Abgleich mit anderen Texten zwecks Auffindung von Übereinstimmungen habe.
\vspace{1em}

\noindent Ich versichere, dass ich die vorliegende Arbeit oder Teile daraus nicht anderweitig als Prüfungsarbeit eingereicht habe.

\vspace{1cm}

\noindent \parbox{20em}{\hrulefill}

\noindent Jad Haddad, Münster, \today